\subsection*{基本定理的其他证明}

在 {[}8{]} 中，F. Schur 证明，在典范坐标中，(15) 的 \(\psi_{ik}\) 满足微分方程

\[
\frac {d}{d t} \left(t \psi_ {i k} (t a)\right) = \delta_ {i k} + \sum_ {j, l} c _ {j l} ^ {k} t a _ {l} \psi_ {i j} (t a).\tag{23}
\]

这些方程经积分后给出一个与下式等价的公式

\[
\varpi (\mathrm{X}) = \sum_ {n > 0} \frac {1}{(n + 1) !} (\mathrm{ad} (\mathrm{X})) ^ {n}\tag{24}
\]

本书第 III 章 §6 第 4 号命题 12 的公式；特别地，在典范坐标中，\(\psi_{ij}\) 可延拓为 \(a_k\) 的整函数。F. Schur 得出一个使 Lie 早先的评论更加精确的结果：若在变换群的定义 (4) 中只假定 \(f_i\) 属于 \(\mathbf{C}^2\) 类，则该群与解析群全同构。 \(\dagger\) E. Cartan 在研究微分系统积分的过程中，于 1904 年引入 Pfaff 形式（{[}12{]}，第 II 卷 \(_{2}\)，第 371 页）

\[
\omega_ {k} = \sum_ {i = 1} ^ {r} \psi_ {k i} d a _ {i} (1 \leqslant i \leqslant r)\tag{25}
\]

（在 (15) 的记号下），后来称为 Maurer--Cartan 形式。Maurer 条件 (17) 可写成

\[
d \omega_ {k} = - \frac {1}{2} \sum_ {i, j} c _ {i j} ^ {k} \omega_ {i} \wedge \omega_ {j};
\]

E. Cartan 证明，有限连续群理论可以从 \(\omega_{k}\) 出发建立，并确立了这一观点与 Lie 观点的等价性。但对他而言，这种方法的意义主要在于它可以适用于“无限连续群”，其理论被他推进得远超 Lie，并使他发展了广义“活动标架”理论。

